\scanpage{128}
of the analogous conjecture already made for relative dimension $1$.

To better grasp this \unsure{second|second} part \inserted{b)} of the conjecture, one would have to examine more closely what becomes, by composition with $\mathcal E_L^{\bar L}\mathrel{\text{---}}\mathcal E_U^{\bar L}$, of the sections of $\mathcal E_L^{\bar L}$ over $\mathcal E_K$ associated with a $V$, and to begin with, study the subgroup \inserted{$\Delta_V$} of $\unsure{\pi|\Pi}_{U/K}^{\bar L}$ image of $I_V$ --- which, one guesses, will play a role analogous to an end-group. To what extent do distinct $V$ define distinct $\Delta_V$?

To study this question, let us introduce a model $X$ of $L/K$ containing $U$, $X$ proper, the irreducible components of \inserted{$D=$} $X-U$ regular divisors \inserted{$D_i$ $(i\in I)$}, and $X-U$ with normal crossings. For $x\in\operatorname{supp}D$, let us consider the henselization \inserted{(not strict)} $\tilde O_x$ of $\underline{O}_{X,x}$, and the inverse image $\bar U_x$ of $U$ over it. We have
\margin{127}
